\chapter{NOTA DELL'AUTORE}\label{nota-dellautore}

\hfill\break

Grazie per aver letto/ascoltato uno dei miei libri!

Ci sono voluti anni per scrivere i miei primi tre titoli: lavoravo come
business manager per una società d'informatica, perciò scrivevo di
notte, nei fine settimana e durante le vacanze. Anche se ho avuto molte
idee per i libri nel corso degli anni, il primo che abbia mai completato
è Assi, e l'ho scritto per le mie nipoti, che al tempo erano
adolescenti. Se leggete Assi, potete vedere alcuni dei primi elementi
delle storie di Expeditionary Force: situazioni impossibili,
problem-solving, pensiero intelligente e un po' di sano sarcasmo.

Poi ho scritto un libro su un programma per sviluppare un tipo di volo
spaziale più veloce della luce: era una storia di avventura riguardo
degli astronauti bloccati su un pianeta alieno che cercavano di
avvertire la Terra di una pericolosa falla nel motore Ftl. Era una buona
storia, e l'ho presentata agli editori tradizionali a metà degli anni
Duemila. E ho ricevuto rifiuti. La mia scrittura era ``solida''; il che,
da allora ho imparato, significa che gli editori non riescono a pensare
a nient'altro da dire, ma non vogliono insultare gli aspiranti
scrittori. La storia però era troppo lunga, volevano che la riducessi a
una novella e cambiassi tutto. Invece di buttare via la storia e
ricominciare, l'ho buttata via e ho provato qualcos'altro.

Il Giorno di Colombo e Ascendente li ho scritti in contemporanea a
partire dal 2011, passavo dalla stesura dell'uno e quella dell'altro.
L'idea di Ascendente mi è venuta dopo avere visto il primo film di Harry
Potter: una delle mie nipoti mi ha chiesto cosa sarebbe successo a Harry
Potter se nessuno gli avesse mai detto che era un mago. Mmh, ho pensato,
questa è un'ottima domanda\ldots{} Così, ho scritto Ascendente.

Nell'originale, primissima versione del Giorno di Colombo, Skippy è un
grazioso, piccolo robot che si è nascosto su una nave quando i kristang
hanno invaso la Terra, e aiuta Joe a sconfiggere gli alieni. Dopo un
anno in cui avevo cercato di scrivere la storia in quella versione,
decisi che sembrava troppo un film della settimana di Disney Channel e,
beh, faceva schifo. Anche se è stato doloroso sprecare un anno di
scrittura, ho buttato via quella versione e ho ricominciato. Questa
volta ho scritto una bozza per l'intero arco della storia di
Expeditionary Force, così avrei saputo dove sarebbe andata a parare. È
stata una grande idea, e mi sono attenuto a quello schema (con alcune
piccole deviazioni lungo la strada).

Con Assi, Il Giorno di Colombo e Ascendente finiti prima dell'estate del
2015, mia moglie mi ha suggerito di:

\hfill\break

provare l'autopubblicazione dei libri su Amazon;

per l'amor di Dio, tacere della mia incapacità di farmi pubblicare i
libri;

pulire il garage.

\hfill\break

Ci sono voluti sei mesi di ricerche e revisioni per preparare i tre
libri da caricare su Amazon. Oltre a riformattare i testi secondo gli
standard richiesti, ho dovuto acquistare le copertine e creare un
account Amazon da scrittore. Quando ho cliccato il pulsante ``Carica2,
il 10 gennaio 2016, la mia più grande speranza era che qualcuno,
chiunque là fuori, comprasse uno dei miei libri, perché allora avrei
potuto essere davvero un autore pubblicato. Dopo aver venduto una copia
di ogni libro, il mio obiettivo era quello di fare abbastanza soldi da
pagare l'immagine di copertina che avevo comprato online (circa
trentacinque dollari per ogni libro).

Per la prima quindicina di gennaio 2016, Amazon ci ha inviato un assegno
di 410,09 dollari, e abbiamo usato parte dei soldi per una bella cena.
Credo che il resto sia andato nell'acquisto di pneumatici nuovi per la
mia auto.

Quando ho caricato Il Giorno di Colombo, ero circa a metà della stesura
del primo capitolo di Operazioni speciali, e continuavo a scrivere di
notte e nei fine settimana. Entro aprile, le vendite del Giorno di
Colombo erano a un numero tale per cui mia moglie e io abbiamo detto:
«Alt, questo potrebbe essere più di un semplice hobby». A quel punto, mi
sono preso una settimana di vacanza per stare a casa a scrivere
Operazioni speciali. Dodici ore al giorno, per nove giorni. Vacanza
davvero divertente! Fare questo ha impresso una bella accelerazione al
lavoro, e Operazioni speciali è stato pubblicato all'inizio di giugno
2016. A metà luglio, con nostro grande stupore, stavamo discutendo
dell'ipotesi di lasciare il mio lavoro per scrivere a tempo pieno. Ad
agosto ho avuto un momento della serie ``la vita è troppo breve'',
quando un amico di famiglia è morto, e poi è mancata anche mia nonna, e
abbiamo deciso di provare questa cosa della scrittura a tempo pieno.
Prima di licenziarmi, ho mostrato a mia moglie un business plan,
elencando i libri che avevo intenzione di scrivere nei tre anni
successivi, con abbozzi di trama e date di pubblicazione. Questo l'ha
rassicurata sul fatto che lasciare il mio lavoro non era una scusa per
stare seduto in pantaloncini e maglietta a guardare film di fantascienza
``per fare ricerca''.

Durante l'estate del 2016, hanno proposto a R.C. Bray di incidere la
narrazione del Giorno di Colombo, e sono sicuro che il suo primo
pensiero è stato: ``Un libro su una lattina di birra parlante? Beeene.
No''. Per fortuna, ci ha ripensato, o ha preso dei farmaci pesanti per
un brutto raffreddore, oppure, se non fosse stato impegnato a registrare
il libro, sua moglie si sarebbe aspettata che ridipingesse casa.
Comunque, R.C. ha registrato Il Giorno di Colombo, è tornato alla sua
favolosa vita, ossia uscire con le star del cinema e colpire palline da
golf facendole uscire dal suo yacht, ed è probabile che si sia
dimenticato della lattina di birra parlante.

Quando ho sentito che R.C. Bray avrebbe narrato Il Giorno di Colombo, la
mia reazione è stata: ``Quell'R.C. Bray? Il tipo che ha narrato I
marziani? Vincitore dell'Audie Award per il miglior narratore di
fantascienza? Ah ah, questa è buona. Ok, chi narrerà davvero il
libro?''.

Poi l'audiolibro del Giorno di Colombo è diventato un grande successo.
Ed è finalista Audie come audiolibro dell'anno!

Quando ho ricevuto l'offerta di creare versioni audio della serie di
Ascendente, mi è stato detto che il narratore sarebbe stato Tim Gerard
Reynolds. La mia reazione è stata: ``Intendi un altro tizio di nome Tim
Gerard Reynolds? Non il Tgr che ha narrato gli audiolibri della serie
Red Rising, giusto?''.

È evidente che ho avuto molta fortuna con i narratori per i miei
audiolibri. Per essere chiari, sono stati loro a scegliere di lavorare
con me, non sono stato io a sceglierli. Se avessi contattato
direttamente Bob o Tim, sarei entrato in modalità super-fanboy e
avrebbero chiesto un ordine restrittivo. Perciò, di nuovo, sono stato
fortunato che abbiano firmato per i progetti.

Finora, non c'è alcun accordo per trarre da Expeditionary Force un film
o uno show televisivo, anche se ho ricevuto richieste da parte di
produttori e studi per i ``diritti di intrattenimento''. Da quello che
la gente del settore mi ha detto, anche se uno studio o una rete
acquisissero i diritti, passerebbe un tempo parecchio lungo prima che
qualsiasi cosa di fatto accadesse. Mi ecciterei per niente e
passerebbero gli anni, mentre il progetto attraverserebbe cicli
interminabili, con produttori e registi che salgono a bordo e scompaiono
e, proprio quando mi sarei completamente arreso, affondato nella fossa
della disperazione, un miracolo accadrebbe e il progetto otterrebbe il
finanziamento! Wow. Non ci conto. D'altra parte, Disney ritirerà i
propri contenuti da Netflix il prossimo anno, quindi Netflix sarà alla
ricerca di nuovi contenuti originali\ldots{}

Ancora una volta, grazie a VOI per aver letto uno dei miei libri. La
scrittura mi fornisce un'ottima scusa per evitare di pulire il garage.